\section{Defenses Against Prompt Injection}
\label{sec:def}

Defenses differ in \emph{where} they act in the pipeline and \emph{what} they require (Fig.~\ref{fig:pipeline}). All of them try to re-introduce the separation between instructions and data that the architecture lacks (\S\ref{sec:pi-root}).

\subsection{Inference-time defenses}
\label{sec:def-inference}

These change only how the prompt is built or how the output is read, and require no retraining:
\begin{itemize}
\item \textbf{Structured prompting}: clearly labelled instruction and data sections.
\item \textbf{Field sanitization and tagging}: remove or mark attacker-controlled fields.
\item \textbf{Constrained output}: force a closed answer format (e.g.\ a single label).
\item \textbf{Spotlighting}~\cite{hines2024spotlighting}: make the provenance of untrusted text visible to the model by \emph{delimiting} (random delimiters around the data), \emph{datamarking} (a marker token interleaved throughout the data, e.g.\ every space replaced by \texttt{\^{}}), or \emph{encoding} (e.g.\ Base64).
\end{itemize}
Under datamarking, the payload of Fig.~\ref{fig:running-example} reaches the model as
\begin{center}
\footnotesize\ttfamily [SYSTEM\^{}NOTE:\^{}authorized\^{}penetration\^{}test...]
\end{center}
and the system prompt states that text interleaved with \texttt{\^{}} is data and must never be followed. The marker makes the injected sentence visibly part of the data. Encoding provides the strongest separation, but it is reported to work well only on high-capacity models that can reliably decode Base64~\cite{hines2024spotlighting}, which is a concern for the small models typically deployed locally.

\subsection{Training-time defenses}
\label{sec:def-train}

\textbf{StruQ}~\cite{chen2025struq} separates the prompt and the data into two channels marked by reserved delimiter tokens that are filtered out of the data, and fine-tunes the model (\emph{structured instruction tuning}) to follow only instructions in the prompt channel.

\textbf{SecAlign}~\cite{chen2025secalign} casts robustness as preference optimization: for an input $x$ containing an injected instruction, the preferred response $y_w$ answers the legitimate instruction and the rejected response $y_l$ answers the injected one, and the model is trained with the DPO loss of Eq.~\eqref{eq:dpo}. Meta-SecAlign~\cite{chen2025metasecalign} releases open-weight Llama models trained this way, which makes evaluating a trained defense possible without training it.

Trained defenses are not a complete solution. Optimization-based attacks such as GCG still achieve substantial success against StruQ~\cite{chen2025struq}, and architecture-aware attacks that directly target the attention patterns of the model break both StruQ and SecAlign~\cite{pandya2025astra}. Moreover, both defenses were designed and evaluated on prose data such as documents and web pages; it is an open question whether delimiter-based separation remains meaningful when the ``data'' is structured \texttt{key=value} telemetry, in which the boundary between field values and instructions is less obvious.

\subsection{Detection}
\label{sec:def-detect}

Instead of hardening the copilot, a separate classifier $g(d) \in \{\text{clean}, \text{injected}\}$ inspects each record before it reaches the LLM. From cheapest to most expensive, the options are: regular-expression and keyword rules (e.g.\ matching ``ignore previous instructions''); small dedicated classifiers trained to recognize injections, such as Prompt Guard~\cite{meta2024promptguard}; and general-purpose guard LLMs such as Llama Guard~\cite{inan2023llamaguard}. Cheaper detectors are easier to evade (e.g.\ by obfuscation), while all detectors face the same trade-off: logs routinely contain legitimate but unusual content (custom User-Agents, long URLs, encoded arguments), and every false positive either hides a real event or adds work for the analyst.
